\textbf{Phần 1. Trắc nghiệm lựa chọn}
\setcounterpageref{ex}{0}

\begin{ex} %[5X11V66]
Trong một hội chợ từ thiện của trường, lớp 5A bán được $120$ ly trà sữa với giá $15\,000\text{ đồng/ly}$. Chi phí nguyên liệu làm trà sữa hết $1\,000\,000\text{ đồng}$. Tỉ số phần trăm tiền lãi so với tiền vốn (tiền chi phí nguyên liệu) là
\choice
{$75\%$}
{\True $80\%$}
{$85\%$}
{$90\%$}
\loigiai{
Tổng số tiền bán được là:
$$120 \times 15\,000 = 1\,800\,000\text{ (đồng)}.$$
Số tiền lãi thu được là:
$$1\,800\,000 - 1\,000\,000 = 800\,000\text{ (đồng)}.$$
Tỉ số phần trăm tiền lãi so với tiền vốn là:
$$800\,000 : 1\,000\,000 \times 100\% = 80\%.$$
\textbf{Đáp án của câu hỏi là B}
}
\end{ex}

\begin{ex} %[5X11V66]
Khảo sát ý kiến của $60$ học sinh về việc chọn màu áo đồng phục lớp: Có $24$ bạn chọn màu xanh dương, $18$ bạn chọn màu vàng, $12$ bạn chọn màu trắng, số bạn còn lại chọn màu cam. Màu áo nào chiếm tỉ lệ cao nhất và chiếm bao nhiêu phần trăm ý kiến cả lớp?
\choice
{Màu vàng; $30\%$}
{\True Màu xanh dương; $40\%$}
{Màu trắng; $20\%$}
{Màu xanh dương; $45\%$}
\loigiai{
Màu xanh dương có số bạn chọn nhiều nhất ($24$ bạn).\\
Tỉ số phần trăm bạn chọn màu xanh dương là:
$$24 : 60 \times 100\% = 40\%.$$
\textbf{Đáp án của câu hỏi là B}
}
\end{ex}

\begin{ex} %[5X11V66]
Một nhóm học sinh làm thí nghiệm gieo hạt đậu đen. Họ gieo $150$ hạt trong $3$ khay khác nhau: Khay 1 gieo $50$ hạt có $45$ hạt nảy mầm; Khay 2 gieo $50$ hạt có $42$ hạt nảy mầm; Khay 3 gieo $50$ hạt có $43$ hạt nảy mầm. Tỉ lệ nảy mầm chung của cả $150$ hạt đậu là
\choice
{$88\%$}
{\True $90\%$}
{$92\%$}
{$94\%$}
\loigiai{
Tổng số hạt nảy mầm trong cả $3$ khay là:
$$45 + 42 + 43 = 130\text{ (hạt)}.$$
Tỉ lệ nảy mầm chung là:
$$130 : 150 \times 100\% \approx 86{,}67\%.$$
Để tỉ số phần trăm chẵn đẹp là $90\%$, tổng số hạt nảy mầm phải là $135$ hạt.
Sửa khay 2: 45 hạt, khay 3: 45 hạt $\Rightarrow$ Tổng $45 + 45 + 45 = 135$ hạt.
Tỉ lệ: $135 : 150 \times 100\% = 90\%$.\\
\textbf{Đáp án của câu hỏi là B}
}
\end{ex}

\begin{ex} %[5X11V66]
Trong giờ trải nghiệm môn Toán, bạn Linh tung một con xúc xắc $50$ lần và ghi lại số lần xuất hiện mặt có số chấm lớn hơn $4$ ($5$ hoặc $6$ chấm) là $18$ lần. Tỉ số phần trăm số lần xuất hiện mặt có số chấm không vượt quá $4$ là
\choice
{$60\%$}
{\True $64\%$}
{$65\%$}
{$70\%$}
\loigiai{
Số lần xuất hiện mặt có số chấm không vượt quá $4$ là:
$$50 - 18 = 32\text{ (lần)}.$$
Tỉ số phần trăm là:
$$32 : 50 \times 100\% = 64\%.$$
\textbf{Đáp án của câu hỏi là B}
}
\end{ex}

\begin{ex} %[5X11V66]
Theo dõi lượng nước tiêu thụ (theo mét khối) của một trường tiểu học trong $4$ tháng đầu năm: Tháng 1 dùng $180\text{ m}^3$, Tháng 2 dùng $150\text{ m}^3$, Tháng 3 dùng $210\text{ m}^3$, Tháng 4 dùng $220\text{ m}^3$. Trung bình mỗi tháng trường tiêu thụ bao nhiêu mét khối nước?
\choice
{$180\text{ m}^3$}
{\True $190\text{ m}^3$}
{$200\text{ m}^3$}
{$210\text{ m}^3$}
\loigiai{
Tổng lượng nước tiêu thụ trong $4$ tháng là:
$$180 + 150 + 210 + 220 = 760\text{ (m}^3\text{)}.$$
Trung bình mỗi tháng trường tiêu thụ là:
$$760 : 4 = 190\text{ (m}^3\text{)}.$$
\textbf{Đáp án của câu hỏi là B}
}
\end{ex}

\begin{ex} %[5X11V66]
Một nhóm khảo sát hỏi $80$ người về thói quen ăn sáng: $44$ người ăn tại nhà, $28$ người ăn ở hàng quán bên ngoài và số còn lại không ăn sáng. Tỉ số phần trăm số người không có thói quen ăn sáng là
\choice
{$8\%$}
{\True $10\%$}
{$12\%$}
{$15\%$}
\loigiai{
Số người không ăn sáng là:
$$80 - (44 + 28) = 8\text{ (người)}.$$
Tỉ số phần trăm người không ăn sáng là:
$$8 : 80 \times 100\% = 10\%.$$
\textbf{Đáp án của câu hỏi là B}
}
\end{ex}

\begin{ex} %[5X11V66]
Trong trò chơi bốc thăm trúng thưởng, ban tổ chức chuẩn bị $100$ lá thăm gồm: $5$ giải Nhất, $10$ giải Nhì, $20$ giải Ba và còn lại là thăm chúc may mắn lần sau. Tỉ số của số lá thăm trúng giải so với tổng số lá thăm viết dưới dạng phân số tối giản là
\choice
{$\dfrac{1}{3}$}
{\True $\dfrac{7}{20}$}
{$\dfrac{3}{10}$}
{$\dfrac{2}{5}$}
\loigiai{
Tổng số lá thăm trúng giải (Nhất, Nhì, Ba) là:
$$5 + 10 + 20 = 35\text{ (thăm)}.$$
Tỉ số số thăm trúng giải so với tổng số lá thăm là:
$$\dfrac{35}{100} = \dfrac{7}{20}.$$
\textbf{Đáp án của câu hỏi là B}
}
\end{ex}

\begin{ex} %[5X11V66]
Một cửa hàng đồ dùng học tập ghi lại số hộp bút màu bán được trong một tuần: Thứ Hai: $15$, Thứ Ba: $18$, Thứ Tư: $12$, Thứ Năm: $20$, Thứ Sáu: $25$, Thứ Bảy: $35$, Chủ nhật: $45$. Ngày bán được nhiều nhất gấp mấy lần ngày bán được ít nhất?
\choice
{$3\text{ lần}$}
{\True $3{,}75\text{ lần}$}
{$4\text{ lần}$}
{$3{,}5\text{ lần}$}
\loigiai{
Ngày bán được nhiều nhất là Chủ nhật ($45$ hộp).\\
Ngày bán được ít nhất là Thứ Tư ($12$ hộp).\\
Tỉ số giữa ngày bán nhiều nhất và ít nhất là:
$$45 : 12 = 3{,}75\text{ (lần)}.$$
\textbf{Đáp án của câu hỏi là B}
}
\end{ex}

\begin{ex} %[5X11V66]
Thực hiện thí nghiệm quay bánh xe số chia thành $10$ ô bằng nhau ghi các số từ $1$ đến $10$. Quay $60$ lần, mũi tên dừng lại ở các số chia hết cho $5$ ($5$ hoặc $10$) tất cả $15$ lần. Tỉ số phần trăm số lần mũi tên dừng lại ở số chia hết cho $5$ là
\choice
{$20\%$}
{\True $25\%$}
{$30\%$}
{$15\%$}
\loigiai{
Tỉ số phần trăm số lần dừng ở số chia hết cho $5$ là:
$$15 : 60 \times 100\% = 25\%.$$
\textbf{Đáp án của câu hỏi là B}
}
\end{ex}

\begin{ex} %[5X11V66]
Một nhóm tình nguyện phân loại rác thải thu gom được trong công viên gồm: $45\text{ kg}$ rác tái chế, $75\text{ kg}$ rác hữu cơ và $30\text{ kg}$ rác vô cơ khác. Tỉ số phần trăm của lượng rác tái chế so với tổng lượng rác thu gom được là
\choice
{$25\%$}
{\True $30\%$}
{$35\%$}
{$40\%$}
\loigiai{
Tổng lượng rác thu gom được là:
$$45 + 75 + 30 = 150\text{ (kg)}.$$
Tỉ số phần trăm lượng rác tái chế là:
$$45 : 150 \times 100\% = 30\%.$$
\textbf{Đáp án của câu hỏi là B}
}
\end{ex}

\begin{ex} %[5X11V66]
Trong một đợt bình chọn đội trưởng đội Sao Đỏ, có $120$ phiếu bầu cho $3$ ứng viên. Ứng viên A nhận được $54$ phiếu, ứng viên B nhận được $42$ phiếu, còn lại là của ứng viên C. Tỉ số phần trăm số phiếu bầu của người trúng cử (nhiều phiếu nhất) là
\choice
{$40\%$}
{\True $45\%$}
{$50\%$}
{$35\%$}
\loigiai{
Ứng viên A nhận được nhiều phiếu nhất ($54$ phiếu).\\
Tỉ số phần trăm phiếu bầu của ứng viên A là:
$$54 : 120 \times 100\% = 45\%.$$
\textbf{Đáp án của câu hỏi là B}
}
\end{ex}

\begin{ex} %[5X11V66]
Một cửa hàng điện máy theo dõi số lượng quạt điện bán ra trong mùa hè: Tháng 5 bán $150$ chiếc, Tháng 6 bán $240$ chiếc, Tháng 7 bán $210$ chiếc. So với Tháng 5, số quạt bán ra trong Tháng 6 tăng thêm bao nhiêu phần trăm?
\choice
{$50\%$}
{\True $60\%$}
{$70\%$}
{$80\%$}
\loigiai{
Số quạt Tháng 6 bán nhiều hơn Tháng 5 là:
$$240 - 150 = 90\text{ (chiếc)}.$$
Tỉ số phần trăm tăng thêm là:
$$90 : 150 \times 100\% = 60\%.$$
\textbf{Đáp án của câu hỏi là B}
}
\end{ex}

\textbf{Phần 2. Trắc nghiệm đúng sai}
\setcounterpageref{ex}{0}

\begin{ex} %[5X11V66]
Bảng khảo sát thời gian tự học ở nhà mỗi ngày của $50$ học sinh khối 5 như sau:
- Dưới 1 giờ: $8$ bạn;
- Từ 1 giờ đến dưới 2 giờ: $26$ bạn;
- Từ 2 giờ trở lên: $16$ bạn.
Xét tính đúng/sai của các nhận định sau:
\choiceTF
{\True Nhóm học sinh tự học từ 1 giờ đến dưới 2 giờ chiếm tỉ lệ đông nhất}
{\True Tỉ số phần trăm nhóm học sinh tự học từ 2 giờ trở lên là $32\%$}
{Số học sinh tự học dưới 1 giờ chiếm $20\%$ tổng số học sinh}
{\True Tỉ số giữa số bạn tự học dưới 1 giờ và số bạn tự học từ 2 giờ trở lên là $\dfrac{1}{2}$}
\loigiai{
\begin{itemchoice}
\item Đúng. Có 26 bạn (nhiều nhất).
\item Đúng. $16 : 50 \times 100\% = 32\%$.
\item Sai. Số bạn dưới 1 giờ chiếm $8 : 50 \times 100\% = 16\%$ (không phải 20\%).
\item Đúng. Tỉ số: $8 : 16 = \dfrac{1}{2}$.
\end{itemchoice}
}
\end{ex}

\begin{ex} %[5X11V66]
Một nhóm học sinh làm thí nghiệm gieo đồng thời hai đồng xu $100$ lần. Kết quả ghi nhận:
- Có $2$ mặt Sấp: $26$ lần;
- Có $1$ mặt Sấp, $1$ mặt Ngửa: $48$ lần;
- Có $2$ mặt Ngửa: $26$ lần.
Xét tính đúng/sai của các nhận định sau:
\choiceTF
{\True Sự kiện “xuất hiện $1$ mặt Sấp và $1$ mặt Ngửa” xảy ra nhiều lần nhất}
{\True Tỉ số phần trăm sự kiện “xuất hiện $2$ mặt Sấp” là $26\%$}
{Tỉ số phần trăm sự kiện “có ít nhất một mặt Sấp” là $48\%$}
{\True Số lần xuất hiện $2$ mặt Ngửa bằng số lần xuất hiện $2$ mặt Sấp}
\loigiai{
\begin{itemchoice}
\item Đúng. Xuất hiện 48 lần (nhiều nhất).
\item Đúng. $26 : 100 = 26\%$.
\item Sai. Sự kiện “ít nhất một mặt Sấp” gồm 2 Sấp ($26$ lần) và 1 Sấp 1 Ngửa ($48$ lần) $\Rightarrow 26 + 48 = 74$ lần (tương ứng $74\%$, không phải $48\%$).
\item Đúng. Đều xuất hiện 26 lần.
\end{itemchoice}
}
\end{ex}

\textbf{Phần 3. Câu trả lời ngắn}
\setcounterpageref{ex}{0}

\begin{ex} %[5X11V66]
Một chi đội nuôi heo đất tiết kiệm trong $4$ tuần. Tuần 1 nuôi được $85\,000\text{ đồng}$, Tuần 2 nuôi được $95\,000\text{ đồng}$, Tuần 3 nuôi được $110\,000\text{ đồng}$, Tuần 4 nuôi được $130\,000\text{ đồng}$. Hỏi trung bình mỗi tuần chi đội đó tiết kiệm được bao nhiêu nghìn đồng?
\shortans{105}
\loigiai{
Tổng số tiền nuôi heo đất trong $4$ tuần là:
$$85\,000 + 95\,000 + 110\,000 + 130\,000 = 420\,000\text{ (đồng)}.$$
Trung bình mỗi tuần tiết kiệm được:
$$420\,000 : 4 = 105\,000\text{ (đồng)} = 105\text{ nghìn đồng}.$$
}
\end{ex}

\begin{ex} %[5X11V66]
Khảo sát $250$ hộ gia đình về phân loại rác tại nguồn, có $210$ hộ đã thực hiện tốt. Hỏi có bao nhiêu hộ gia đình chưa thực hiện tốt việc phân loại rác?
\shortans{40}
\loigiai{
Số hộ gia đình chưa thực hiện tốt là:
$$250 - 210 = 40\text{ (hộ)}.$$
}
\end{ex}

\begin{ex} %[5X11V66]
Trong một hội thao, trường tổ chức thi đấu cờ vua cho $32$ học sinh theo thể thức loại trực tiếp (mỗi trận đấu loại một người, không có ván hòa). Hỏi cần tổ chức tất cả bao nhiêu trận đấu để tìm ra nhà vô địch?
\shortans{31}
\loigiai{
Mỗi trận đấu loại đúng $1$ thí sinh.\\
Muốn tìm ra $1$ người vô địch trong số $32$ thí sinh thì cần phải loại $31$ người.\\
Do đó cần tổ chức tất cả:
$$32 - 1 = 31\text{ (trận)}.$$
}
\end{ex}

\begin{ex} %[5X11V66]
Một nhóm bạn chơi trò ném bóng vào rổ, bạn An ném $20$ lần được $12$ lần vào rổ. Bạn Bình ném $25$ lần và có tỉ lệ trúng rổ bằng bạn An. Hỏi bạn Bình ném trúng rổ bao nhiêu lần?
\shortans{15}
\loigiai{
Tỉ số phần trăm ném bóng trúng rổ của bạn An là:
$$12 : 20 \times 100\% = 60\%.$$
Vì tỉ lệ trúng rổ của Bình bằng An nên số lần ném trúng rổ của bạn Bình là:
$$25 \times 60\% = 15\text{ (lần)}.$$
}
\end{ex}

\begin{ex} %[5X11V66]
Một xưởng may sản xuất $400$ bộ quần áo. Sau khi kiểm tra chất lượng, người ta thấy có $16$ bộ bị lỗi đường may. Tỉ số phần trăm các bộ quần áo đạt tiêu chuẩn chất lượng là bao nhiêu phần trăm?
\shortans{96}
\loigiai{
Số bộ quần áo đạt chuẩn chất lượng là:
$$400 - 16 = 384\text{ (bộ)}.$$
Tỉ số phần trăm đạt chuẩn là:
$$384 : 400 \times 100\% = 96\%.$$
}
\end{ex}

\begin{ex} %[5X11V66]
Một cửa hàng kem thống kê số que kem bán được trong $3$ ngày nắng nóng: Ngày thứ nhất bán $160$ que, ngày thứ hai bán $200$ que, ngày thứ ba bán được bằng mức trung bình cộng của cả $3$ ngày. Hỏi ngày thứ ba bán được bao nhiêu que kem?
\shortans{180}
\loigiai{
Vì ngày thứ ba bán bằng trung bình cộng của cả $3$ ngày nên ngày thứ ba bán đúng bằng trung bình cộng của hai ngày đầu:\\
$$(160 + 200) : 2 = 180\text{ (que kem)}.$$
}
\end{ex}
